\section{User interfaces: GUI, CLI, and visualizer}\label{sec:interfaces}

QtRocket ships three executables over the one simulation core surveyed in \cref{sec:overview}: the
Qt6 Widgets application \code{qtrocket}, the headless REPL \code{qtrocket-cli}, and the standalone
OpenGL viewer \code{qtrocket-visualizer}. Each links \code{qtrocket\_core} and nothing of the other
two; the visualizer's build file states outright that it does not touch the GUI or CLI code
(\src{visualizer/CMakeLists.txt}{5}). This section shows the document's recurring motif in its most
visible form: capability accumulates in the core, and the front ends expose sharply uneven subsets
of it. The CLI is the broadest surface --- it can author, persist, and fly a multi-part design; the
visualizer is the most finished piece of visual UI in the project; and the GUI, the executable that
carries the project's name, is the narrowest --- a motor-selection-and-launch panel whose
design-tool half is empty scaffolding at commit \code{254cd41}.

\subsection{The Qt6 GUI: a motor-and-launch front end}\label{sec:interfaces-gui}

\code{main.cpp} instantiates the logger and the \code{QtRocket} controller singleton, then blocks in
\cppinline|gui::run| (\src{main.cpp}{25}). That function --- the only bridge between the Qt-free
controller and Qt --- spawns a \cppinline|std::thread| running \code{guiWorker}, which constructs
\code{QApplication}, installs a translator, creates \code{MainWindow}, and enters
\cppinline|exec()| (\src{gui/GuiRunner.cpp}{30}), then immediately joins it.

\begin{hardfail}[QApplication on a spawned thread]
\cppinline|std::thread guiThread(guiWorker, ...); guiThread.join();| --- the Qt event loop runs on
a secondary thread while the main thread blocks in \cppinline|join()| for the application's
lifetime (\src{gui/GuiRunner.cpp}{59}). The indirection buys nothing --- the caller does no
concurrent work --- while putting the GUI on a non-main thread, which Qt only tolerates on some
platforms and which is a hard failure on macOS, where the GUI must run on the first thread. \fail
\end{hardfail}

Startup also performs two loads that can never succeed. The application icon is read from
\code{:/qtrocket.png} (\src{gui/GuiRunner.cpp}{31}), a path absent from the project's only
\code{.qrc}, whose sole entry is \code{resources/rocket64.png} under \code{/images}
(\src{qtrocket.qrc}{4}); the icon silently resolves to null, masked because the main window loads
the valid path separately (\src{gui/MainWindow.ui}{16}). And the internationalization scaffolding
is dead code: the translator loads \code{qtrocket\_<locale>} from a \code{:/i18n/} resource prefix
that is never populated (\src{gui/GuiRunner.cpp}{40}) --- the build-side half of the story, translation files
generated but attached to nothing, is covered in \cref{sec:testing} --- and the adjacent \code{TODO} admits the
application is en\_US only (\src{gui/GuiRunner.cpp}{34}).

The main window's layout promises a design tool it does not contain. The \code{.ui} places a
\code{RocketTreeView} (a promoted \code{QTreeView}), a tab widget, and a \code{RocketModelerView}
in the central splitter (\src{gui/MainWindow.ui}{49}); the constructor adds two tabs
programmatically --- \code{CannonballTab} and \code{SimOptionsTab} (\src{gui/MainWindow.cpp}{30})
--- and wires exactly three menu actions: Quit, Save Motor Database, and About
(\src{gui/MainWindow.cpp}{42}). Both design widgets are default-constructed shells with no model,
no data, and no behavior; the complete implementation of the tree widget at \code{HEAD} is
(\src{gui/RocketTreeView.cpp}{1}):

\begin{minted}{cpp}
#include "RocketTreeView.h"

RocketTreeView::RocketTreeView(QWidget* parent) : QTreeView(parent)
{

}
\end{minted}

\noindent \code{RocketModelerView} is the same empty \code{QWidget} shell
(\src{gui/RocketModelerView.cpp}{3}), so the left panel and bottom panel render permanently blank.
The File menu completes the picture: New, Open, Save, Close, and Save~As are declared but ship
\code{enabled=false} in the \code{.ui} --- \code{actionNew} at \src{gui/MainWindow.ui}{118}, Open,
Save, Close, and Save~As at lines 130, 142, 154, and 171 --- with no slots at \code{HEAD}.
Every one of these capabilities exists one layer down: the CLI's \code{newdesign}
(\src{cli/Repl.cpp}{784}) and \code{addpart} (\src{cli/Repl.cpp}{825}) author part trees the GUI
cannot even display. Structural design status in the GUI at the reviewed commit: \fabsent{}
(\fwip{} on the in-flight branch of \cref{sec:interfaces-branch}).

What the GUI does do is launch point masses, and the tab name is honest about it. The
\code{CannonballTab} takes initial velocity, a launch angle validated to $0$--$90^{\circ}$ from
vertical (\src{gui/CannonballTab.cpp}{36}), mass, drag coefficient, and reference area
(\src{gui/CannonballTab.cpp}{40}). The launch slot decomposes the velocity into downrange and
vertical components (\src{gui/CannonballTab.cpp}{104}), writes mass, $C_d$, and area into the shared
\code{RocketModel} through the \code{QtRocket} facade --- the GUI's only route to the core
(\src{gui/CannonballTab.cpp}{110}) --- and calls \cppinline|launchRocket()|
(\src{gui/CannonballTab.cpp}{115}). That call is fully synchronous: it runs
\cppinline|runUntilTerminate()| on the calling thread (\src{QtRocket.cpp}{54}), freezing the window
for the duration of the propagation before \code{AnalysisWindow} opens
(\src{gui/CannonballTab.cpp}{118}). One rule gates the workflow: Calculate Trajectory is enabled
iff the rocket has a motor (\src{gui/CannonballTab.cpp}{88}), a condition reachable by four
acquisition paths: (i) local file import of RockSim \code{.rse} or RASP \code{.eng}, dispatched by
extension with exceptions caught into a critical dialog (\src{gui/CannonballTab.cpp}{136},
\src{gui/CannonballTab.cpp}{141}); (ii) a \code{.qmd} motor-database round trip with suffix
enforcement (\src{gui/CannonballTab.cpp}{173}, \src{gui/CannonballTab.cpp}{212}), duplicated in
the Tools menu (\src{gui/MainWindow.cpp}{74}); (iii) online search through the
\code{ThrustCurveMotorSelector} dialog (\src{gui/CannonballTab.cpp}{163}); and (iv) an engine
combo rebuilt from the database as the single source of truth (\src{gui/CannonballTab.cpp}{156})
plus Set Motor, which resolves the name and installs the motor (\src{gui/CannonballTab.cpp}{231}).

The online selector is a three-click manual workflow --- ``Get Metadata'' populates search facets
(\src{gui/ThrustCurveMotorSelector.cpp}{52}), ``Search'' merges results into the local database
(\src{gui/ThrustCurveMotorSelector.cpp}{75}), ``setMotor'' installs the selection and plots its
thrust curve (\src{gui/ThrustCurveMotorSelector.cpp}{83}) --- and every click performs network I/O
synchronously on the GUI thread; the core documents the cost itself, noting downloads are
``performed synchronously on the caller's (often GUI) thread''
(\src{model/ThrustCurveClient.cpp}{292}), with searches capped at 100 results
(\src{model/ThrustCurveClient.cpp}{295}). Failures are silent by construction:
\code{CurlConnection} returns an empty string on any transport or HTTP failure
(\src{utils/CurlConnection.h}{28}), the client degrades to empty results with Logger-only errors
(\src{model/ThrustCurveClient.cpp}{271}), and the GUI slots have no failure branch
(\src{gui/ThrustCurveMotorSelector.cpp}{48}) --- a user without connectivity clicks Get Metadata
and simply sees nothing happen.

The one polished corner is \code{SimOptionsTab}: a live-apply form whose atmosphere, gravity, and
integrator combos are enumerated from the sim layer's own model registries via throwaway
\code{Environment}/\code{Integrator} instances (\src{gui/SimOptionsTab.cpp}{32}; registries at
\src{sim/Environment.h}{42}, \src{sim/Integrator.h}{43}), so the GUI cannot drift from the core's
vocabulary. A one-time apply syncs \code{QtRocket} to the displayed defaults
(\src{gui/SimOptionsTab.cpp}{71}), and gravity/atmosphere changes mutate the shared
\code{Environment} in place so the two selections do not clobber each other
(\src{gui/SimOptionsTab.cpp}{93}). \fpresent

Post-run analysis is thinner. \code{AnalysisWindow} embeds a vendored QCustomPlot v2.1.1
(\src{gui/qcustomplot.h}{23}) and offers exactly three graphs: altitude
(\cppinline|position[2]| vs.\ time, \src{gui/AnalysisWindow.cpp}{49}), vertical velocity
(\src{gui/AnalysisWindow.cpp}{74}), and the motor thrust curve (\src{gui/AnalysisWindow.cpp}{86}).
The axis-range code dereferences \cppinline|std::min_element| without an emptiness check
(\src{gui/AnalysisWindow.cpp}{111}; repeated at \src{gui/ThrustCurveMotorSelector.cpp}{113}) ---
undefined behavior on an empty series, and an empty default \code{ThrustCurve} does exist
(\src{model/ThrustCurve.h}{17}); only the motor-gated launch flow keeps the path unreachable. More
telling is what the window never shows: the controller exposes whole-flight
\code{TrajectoryStatistics} (\src{QtRocket.h}{55}) and a \code{TerminationReason}
(\src{QtRocket.h}{62}), and the propagator records mass, CG, and inertia alongside every state ---
all of which the CLI surfaces (\src{cli/Repl.cpp}{703}, \src{cli/Repl.cpp}{104}) --- yet neither
accessor has a single call site under \code{gui/}. A user whose run aborted for no-liftoff gets an
altitude plot and no explanation, and must eyeball the curve for apogee. The Barrowman
center-of-pressure machinery of \cref{sec:aero} is likewise computed and tested in the core with no
display surface in any front end. \fpartial

The remaining gaps are individually small and collectively diagnostic: no \code{QSettings} or other
preference persistence exists anywhere in \code{gui/}, \code{visualizer/}, or \code{main.cpp} ---
file dialogs hardcode \code{/home} as their start directory (\src{gui/CannonballTab.cpp}{127},
\src{gui/MainWindow.cpp}{78}) --- dialogs keep the Designer default title ``Dialog''
(\src{gui/AnalysisWindow.ui}{14}), and \code{AnalysisWindow.ui} declares a
\cppinline|plotAltitude()| slot that does not exist in the class (\src{gui/AnalysisWindow.ui}{77}).
None is hard to fix; together they mark a front end that has not yet had a first polish pass.

\subsection{The command-line REPL}\label{sec:interfaces-cli}

\code{qtrocket-cli} is the inverse of the GUI: no widgets, no polish debt, and the widest command
surface in the project. Its \cppinline|main| deliberately ignores \code{argc}/\code{argv}
(\src{cli/CliMain.cpp}{13}) --- no flags, no script-file argument, no \code{--version} --- and sets
the logger to \code{ERROR} because ``Logger shares stdout with the CLI's machine-readable output''
(\src{cli/CliMain.cpp}{15}), then runs \code{cli::Repl} over
\code{std::cin}/\code{std::cout} (\src{cli/CliMain.cpp}{25}). The loop is a bare
\cppinline|std::getline| with a flush per command (\src{cli/Repl.cpp}{232}); there is no readline,
no history, and no prompt --- the project's own \code{CLAUDE.md} calls it a ``readline-style REPL,''
but no such library is linked (\src{CMakeLists.txt}{297}). That austerity is the documented intent:
the class comment promises the same operations the GUI performs, ``interactively, from a pipe, or
from a redirected script'' (\src{cli/Repl.h}{20}), and batch use is plain shell redirection.

\begin{designrule}[The stdout protocol is the interface contract]
Every command answers exactly one \code{OK} or \code{ERR} line (dispatch chain at
\src{cli/Repl.cpp}{245}, unknown-command fallback at \src{cli/Repl.cpp}{981});
\code{launch} appends \code{WARN:} on non-nominal termination --- through an exhaustive switch with
no default, so a new \code{TerminationReason} fails to compile until the CLI handles it
(\src{cli/Repl.cpp}{41}) --- and a \code{CSV <abs-path>} line; informational output is
\code{\#}-prefixed and \code{\#} lines are comments on input (\src{cli/Repl.cpp}{256}). Because the
REPL is built as a static \code{cli} library (\src{CMakeLists.txt}{235}), the test binaries drive
the exact object code users type into and parse this same protocol (\cref{sec:testing}).
\end{designrule}

\Cref{tab:interfaces-repl} inventories the command set. The count deserves precision: the dispatch
chain has 32 handler branches covering 33 command words (\code{quit} and \code{exit} share one).
The grammar for design authoring is \code{key=value} over the part factory
(\src{cli/Repl.cpp}{169}; \src{model/parts/PartFactory.cpp}{27}), with strict full-string numeric
parsing that rejects trailing garbage (\src{cli/Repl.cpp}{119}) and warn-and-skip tolerance for
unknown keys (\src{cli/Repl.cpp}{205}).

\begin{table}[htbp]
\centering\footnotesize
\caption{The \code{qtrocket-cli} command set at commit \code{254cd41}: 33 command words, 32 handler
branches in \code{Repl::execute} (\srccap{cli/Repl.cpp}{245}).}
\label{tab:interfaces-repl}
\begin{tabularx}{\linewidth}{@{}>{\ttfamily}p{0.34\linewidth}X@{}}
\toprule
{\normalfont\textbf{Command}} & \textbf{Behavior} \\
\midrule
\multicolumn{2}{@{}l}{\itshape Session}\\
help & command list, \code{\#}-prefixed \\
quit / exit & ``\code{OK bye}''; ends the loop \\
\addlinespace
\multicolumn{2}{@{}l}{\itshape Motor database}\\
loadmotors <file> & import RockSim \code{.rse} / RASP \code{.eng}, by extension \\
savedb / loaddb <file.qmd> & motor-database XML save; additive load \\
listmotors [substr] & filtered summaries (avg thrust, total impulse) \\
setmotor <code> & resolve by common name, install on the rocket \\
tcfacets & thrustcurve.org search facets (online) \\
tcsearch <k=v>\ldots & online search (manufacturer, diameter, impulse class); merges into the database \\
\addlinespace
\multicolumn{2}{@{}l}{\itshape Simulation configuration}\\
setmass / setdrag / setarea & written through to the \code{RocketModel} immediately \\
setvelocity / setangle & staged; applied at launch \\
settimestep <s> & integrator step, $> 0$ \\
listatmospheres / setatmosphere & names validated against the sim registries; likewise
  \code{listgravity}/\code{setgravity} and \code{listintegrators}/\code{setintegrator} \\
status & staged configuration plus database size \\
\addlinespace
\multicolumn{2}{@{}l}{\itshape Design authoring}\\
newdesign <type> [k=v\ldots] & root part via the \code{makePart} factory \\
addpart <parent|root> <type> [k=v\ldots] & attach with \code{abut(z)}; \code{x}/\code{y} parsed but ignored \\
listparts & indented id/type/name/mass tree; composite mass/CG footer \\
removepart <id> & refuses the root; re-checks motor presence \\
cleardesign & restore the boot placeholder \\
\addlinespace
\multicolumn{2}{@{}l}{\itshape Flight and persistence}\\
launch & run to termination; stats, \code{WARN} on abort, auto-CSV \\
states [stride] & state series as CSV to stdout \\
save <path.csv> & full series to file; echoes the absolute path \\
savedesign / loaddesign <file.qrd> & design round trip (\cref{sec:persistence}) \\
\bottomrule
\end{tabularx}
\end{table}

This inventory is not a paper capability; the review exercised it end to end against the built
binary. A single piped session loaded 249 motors from the bundled \code{Aerotech.rse}
(\src{cli/Repl.cpp}{314}), loaded the committed \code{large38\_multi.qrd} fixture with its motor
re-resolved by name (``\code{OK loaddesign: ... (motor I280DM)}''), listed the part tree with its
composite mass/CG footer (\src{cli/Repl.cpp}{880}), and flew it:
\code{OK launch: 4258 steps, t\_final=42.570s}, apogee $1856.613\unit{m}$ at $t=13.240\unit{s}$,
maximum speed $562.928\unit{m/s}$ at $t=1.730\unit{s}$, statistics read from the live-tracked
\code{TrajectoryStatistics} rather than re-derived (\src{cli/Repl.cpp}{703}). A second fixture
under the adaptive Runge-Kutta-Fehlberg integrator completed in 440 steps, an order of magnitude
fewer than the fixed-step run, and a CLI-built design round-tripped through
\code{savedesign}/\code{cleardesign}/\code{loaddesign} with identical part masses
(\cref{sec:persistence} covers the mechanics).

The REPL's gaps are edges, not holes. The process exit code is always zero
(\src{cli/Repl.cpp}{242}; the contract comment says so, \src{cli/Repl.h}{28}), so the scripting
mode that is the CLI's stated strength cannot signal failure to a shell --- callers must parse
\code{ERR} lines. Every \code{launch} unconditionally writes \code{qtrocket\_run.csv} into the
current directory under a hardcoded name (\src{cli/Repl.cpp}{716}). \code{addpart} parses
\code{x=} and \code{y=} into an offset vector (\src{cli/Repl.cpp}{189}) and silently discards
them, attaching with \cppinline|abut(offset.z())| only (\src{cli/Repl.cpp}{872}) --- and that abut
is the only seat kind reachable from any front end, although the \code{.qrd} format round-trips
all three (\cref{sec:persistence}). No command sets the design name, so although the serializer
writes and restores it (\src{model/DesignSerializer.cpp}{221}), every CLI-authored file carries
\code{<design name=""/>} (\src{cli/Repl.cpp}{259}). \code{status} prints a staged
\code{dry\_mass} that \code{loaddesign} never re-syncs (\src{cli/Repl.cpp}{966}): after loading a
$0.785\unit{kg}$ design, the live session still reported ``\code{dry\_mass = 1 kg}''
(\src{cli/Repl.cpp}{656}) --- harmless to the physics, which uses the composite tree mass, but
misleading in the one place a user checks configuration. Two smaller notes: \code{DEG\_PER\_RAD}
is a truncated copy of the GUI's constant ($57.2958$ vs.\ $180/\pi = 57.29578\ldots$) whose
comment points at the wrong file (\src{cli/Repl.cpp}{35}); and \code{setangle} tilts only the
initial velocity vector (\src{cli/Repl.cpp}{679}) --- thrust remains world-vertical throughout the
burn, a 3-DOF core limitation owned by \cref{sec:sim}.

\subsection{The standalone visualizer}\label{sec:interfaces-viz}

\code{qtrocket-visualizer} is built unconditionally from the root build
(\src{CMakeLists.txt}{243}) and links only the Qt OpenGL modules plus \code{qtrocket\_core}
(\src{visualizer/CMakeLists.txt}{27}). Its \cppinline|main| disables the RHI widget backing store
and deliberately leaves the GL version and profile unconstrained --- the documented rationale is
surviving EGL/Wayland/NVIDIA context matching (\src{visualizer/main.cpp}{10}); an optional
\code{.qrd} path on \code{argv} is opened at startup (\src{visualizer/main.cpp}{35}). The window
owns a headless \code{RocketModel} and a deliberately empty \code{MotorModelDatabase} purely to
satisfy the loader's signature (\src{visualizer/VisualizerWindow.h}{69}) --- a missing motor is a
harmless warning, since only geometry matters here --- and opens designs through the same
\code{model::DesignSerializer::load} the CLI uses (\src{visualizer/VisualizerWindow.cpp}{184}),
keeping the current model on failure (\src{visualizer/VisualizerWindow.cpp}{202}).

The architectural keystone is where the meshes get their positions
(\src{visualizer/RocketMesh.cpp}{429}):

\begin{minted}{cpp}
   // Consume the same absolute placement the simulator does: resolve every part's pose once, then
   // translate each origin-centered primitive so its fore plane lands at the resolved fore-plane
   // origin (the build* primitives are centered about mid-length, half an axialLength forward of
   // the aft plane).
   const std::vector<model::part::Placed> placed =
      model::part::resolvePlacements(*root, model::part::Pose{});
\end{minted}

\begin{keyinsight}[One resolver, two consumers]
\code{buildRocketMeshes} positions every primitive from
\cppinline|model::part::resolvePlacements| --- the same routine the simulator's composite
mass/CM/inertia cache consumes (\src{visualizer/RocketMesh.cpp}{433};
\src{model/parts/Part.cpp}{166}) --- so rendered geometry and simulated geometry cannot diverge by
construction. This closes, for the visualizer, exactly the silent renderer-vs-simulator divergence
that motivated the placement redesign of \cref{sec:model}.
\end{keyinsight}

The same walk reads the cached overlap-diagnostics verdict from the part tree
(\src{visualizer/RocketMesh.cpp}{440}; \src{model/parts/Part.h}{128}) and tags offending parts,
which the GL widget paints in a fixed error red overriding any color scheme
(\src{visualizer/RocketGLWidget.cpp}{621}) --- the physics layer's self-intersection check made
visible. Characteristically for this codebase, the diagnostic message is carried on each render
item ``for a tooltip / status line'' that no UI consumes yet (\src{visualizer/RocketMesh.h}{53}).

Mesh coverage is dispatched by \cppinline|dynamic_cast| over the concrete part types
(\src{visualizer/RocketMesh.cpp}{456}): a cone for \code{ConicalNoseCone}, hollow walls with
annular caps for \code{BodyTube} (\src{visualizer/RocketMesh.cpp}{153}), flat-plate trapezoidal
prisms arrayed around the body for \code{FinSet} (\src{visualizer/RocketMesh.cpp}{335}), and a UV
sphere for \code{HollowSphere}; \code{Motor} and unknown types get no geometry
(\src{visualizer/RocketMesh.cpp}{478}), so a loaded design's motor is invisible. Primitives are
origin-centered with $+z$ toward the nose (\src{visualizer/RocketMesh.h}{70}) and translated so the
fore plane lands at the resolved origin (\src{visualizer/RocketMesh.cpp}{482}). The rendering layer
is genuinely portable: dual shader dialects (GLSL 1.20 desktop, GLSL ES 1.00) chosen via
\cppinline|context()->isOpenGLES()| (\src{visualizer/RocketGLWidget.cpp}{314}), wireframe as a
pre-expanded \code{GL\_LINES} edge buffer because \cppinline|glPolygonMode| is unavailable on
GLES2 (\src{visualizer/RocketGLWidget.cpp}{518}), and an orbit/pan/zoom camera with a
bounding-sphere reset (\src{visualizer/RocketGLWidget.cpp}{434},
\src{visualizer/RocketGLWidget.cpp}{261}). Four editable color-scheme presets map
\cppinline|Part::typeName()| to colors (\src{visualizer/ColorScheme.cpp}{22}), but the per-type
swatch panel is hardcoded to exactly the four current type names
(\src{visualizer/VisualizerWindow.cpp}{44}); a fifth part type would render in the fallback color
with no picker.

The one thing the visualizer lacks is a caller. It is not integrated with the GUI in any way ---
the GUI never launches or embeds it (\src{gui/MainWindow.cpp}{22}) --- and at \code{HEAD} the GUI
cannot even produce the \code{.qrd} files it consumes; only the CLI's \code{savedesign} can. The
most finished visual component in the project is unreachable from the product's main window.
\fpresent{} as a tool, \fabsent{} as an integrated feature.

\subsection{Work in flight: the
  \texorpdfstring{\code{RocketTreeViewImplementation}}{RocketTreeViewImplementation} branch}%
\label{sec:interfaces-branch}

Everything above describes commit \code{254cd41}. An in-flight branch,
\code{RocketTreeViewImplementation}, sits two commits ahead (\code{049bb02} ``Initial
implementation of RocketTreeView Widget'', \code{2ddd553} ``GUI save/load WIP''; $+360/-15$ lines
across ten files) and closes part of the GUI's design gap. It adds \code{RocketPartModel}, a
read-only \code{QAbstractItemModel} over the part tree with Name/Type/Mass columns, using raw
\code{Part*} internal pointers and a new \code{Part::getParent()} accessor for the parent walk
(declared at \src{model/parts/Part.h}{212} on that branch). The tree stays live through a
deliberately Qt-free seam: \code{RocketModel::setStructureChangedCallback}, a
\cppinline|std::function| observer hook added at \src{model/RocketModel.h}{121} on the branch and
fired from \code{setMotorModel}, \code{setRoot}, \code{addPart}, and \code{removePart}; the view
re-roots and expands on every fire and detaches in its destructor. The branch also wires
File\,>\,Open/Save/Save~As to \code{DesignSerializer} for \code{.qrd} files with current-file
tracking and window-title updates (handlers in the region of \src{gui/MainWindow.cpp}{92} on the
branch), and adds \code{CannonballTab::refreshFromModel()} so a load re-syncs the mass/$C_d$/area
fields rather than letting the next Calculate silently overwrite loaded values. Even on the branch
the tree is display-only --- no part can be added, removed, or edited from the GUI; New and Close
remain disabled; and the refresh is a full model reset per change. The second commit's own title,
``GUI save/load WIP,'' states the status accurately. \fwip

One cross-cutting fact frames all three front ends: nothing under \code{gui/} or
\code{visualizer/} has any automated test (\cref{sec:testing}), so every ``implemented'' verdict in
this section rests on code reading and the live CLI sessions above, not on a suite. And the CLI's
class comment --- ``the same operations the GUI performs \ldots{} as text commands''
(\src{cli/Repl.h}{20}) --- was written when the GUI was the reference point. At the reviewed commit
the relationship has inverted: the sentence now undersells a REPL that authors, persists, and flies
designs the GUI cannot yet open. The consolidated ranking of these gaps belongs to
\cref{sec:roadmap}.
